% GENERATED FILE. Edit canonical lessons, then run npm run generate:books.
% canonical-source-hash: fnv1a64:26928a221c5c65aa
% canonical-lessons: TA-C509-patavi-uyarvu, TA-C509-latciyam, TA-C509-ulai, TA-C509-ukkuvi, TA-C509-nipunar

\chapter{Climbing the Ladder}
\label{ch:ta-climbing-the-ladder}
\hlchaptermodality{509}

\begin{chapteropening}
\textbf{By the end of this chapter:} \emph{I can talk about promotions and ambitions.}

Complete the last lesson of chapter 509: I can talk about promotions and ambitions.
\end{chapteropening}

\section[\texorpdfstring{\ta{பதவி} \ta{உயர்வு} (patavi uyarvu)}{பதவி உயர்வு (patavi uyarvu)}]{\ta{பதவி} \ta{உயர்வு} (patavi uyarvu) — a promotion}

\label{lesson:TA-C509-patavi-uyarvu}



\begin{quote}
Before the new one: say the Tamil for an accountant, then the Tamil for a clerk.
\end{quote}

\subsection*{You'll want to know: \ta{பதவி} \ta{உயர்வு}}
\textbf{\ta{பதவி} \ta{உயர்வு}} — \emph{patavi uyarvu} — \textquotedblleft{}a promotion\textquotedblright{}.

\begin{quote}
\ta{கடினமாக} \ta{உழைத்ததால்} \ta{அவளுக்குப்} \ta{பதவி} \ta{உயர்வு} \ta{கிடைத்தது}.
\end{quote}

\begin{quote}
\emph{kaṭiṉamāka uḻaittatāl avaḷukkup patavi uyarvu kiṭaittatu.}
\end{quote}

\begin{quote}
She got a promotion because she worked hard.
\end{quote}

Word for word \textquotedblleft{}a rise in post\textquotedblright{}: \ta{பதவி} \textquotedblleft{}a post\textquotedblright{} + \ta{உயர்வு} \textquotedblleft{}a rise\textquotedblright{}. The person takes -\ta{க்கு}: \ta{அவளுக்குப்} \ta{பதவி} \ta{உயர்வு}.

\begin{cousinweb}
\textbf{\ta{பதவி} \ta{உயர்வு}} (\emph{patavi uyarvu}) joins \ta{பதவி} (\emph{patavi}), \textquotedblleft{}a post\textquotedblright{}, borrowed from Sanskrit \emph{padavī}, \textquotedblleft{}a path, a station\textquotedblright{} (Monier-Williams), and the native \ta{உயர்வு} (\emph{uyarvu}), \textquotedblleft{}a rise\textquotedblright{}, from \ta{உயர்} (\emph{uyar}), \textquotedblleft{}to rise\textquotedblright{} (DEDR 646).
\end{cousinweb}

\begin{grammarlens}[title={What you've built}]
One more word for describing what you lived and felt.
\end{grammarlens}

\subsection*{Your turn}
\begin{itemize}
\raggedright
  \item \emph{Say it:} \emph{patavi uyarvu}
  \item \emph{Say it:} \emph{patavi uyarvu}, once more
  \item \emph{Say it:} \emph{kaṭiṉamāka uḻaittatāl avaḷukkup patavi uyarvu kiṭaittatu.} — She got a promotion because she worked hard.
  \item \emph{Recall:} say the Tamil for an accountant, then the Tamil for a clerk, then say \emph{patavi uyarvu} again
\end{itemize}

\begin{tcolorbox}[breakable,colback=teal!4,colframe=teal!35!black,title={Before you move on}]
What does \ta{பதவி} \ta{உயர்வு} mean? (\textquotedblleft{}A promotion\textquotedblright{}.) Say it once more.
\end{tcolorbox}
\section[\texorpdfstring{\ta{லட்சியம்} (laṭciyam)}{லட்சியம் (laṭciyam)}]{\ta{லட்சியம்} (laṭciyam) — an ambition}

\label{lesson:TA-C509-latciyam}



\begin{quote}
Before the new one: say the Tamil for a clerk, then the Tamil for a promotion.
\end{quote}

\subsection*{You'll want to know: \ta{லட்சியம்}}
\textbf{\ta{லட்சியம்}} — \emph{laṭciyam} — \textquotedblleft{}an ambition\textquotedblright{}.

\begin{quote}
\ta{மருத்துவர்} \ta{ஆவது} \ta{என்} \ta{தம்பியின்} \ta{லட்சியம்}.
\end{quote}

\begin{quote}
\emph{maruttuvar āvatu eṉ tampiyiṉ laṭciyam.}
\end{quote}

\begin{quote}
Becoming a doctor is my younger brother's ambition.
\end{quote}

It is also written \ta{இலட்சியம்}. The plural is \ta{லட்சியங்கள்}.

\begin{cousinweb}
\textbf{\ta{லட்சியம்}} (\emph{laṭciyam}) is borrowed from Sanskrit \emph{lakṣya}, \textquotedblleft{}a mark, something aimed at\textquotedblright{} (Monier-Williams). Tamil writes the Sanskrit \emph{kṣ} as \ta{ட்ச} and gives the noun -\ta{அம்}.
\end{cousinweb}

\begin{grammarlens}[title={What you've built}]
One more word for describing what you lived and felt.
\end{grammarlens}

\subsection*{Your turn}
\begin{itemize}
\raggedright
  \item \emph{Say it:} \emph{laṭciyam}
  \item \emph{Say it:} \emph{laṭciyam}, once more
  \item \emph{Say it:} \emph{maruttuvar āvatu eṉ tampiyiṉ laṭciyam.} — Becoming a doctor is my younger brother's ambition.
  \item \emph{Recall:} say the Tamil for a clerk, then the Tamil for a promotion, then say \emph{laṭciyam} again
\end{itemize}

\begin{tcolorbox}[breakable,colback=teal!4,colframe=teal!35!black,title={Before you move on}]
What does \ta{லட்சியம்} mean? (\textquotedblleft{}An ambition\textquotedblright{}.) Say it once more.
\end{tcolorbox}
\section[\texorpdfstring{\ta{உழை} (uḻai)}{உழை (uḻai)}]{\ta{உழை} (uḻai) — to work hard; to toil}

\label{lesson:TA-C509-ulai}



\begin{quote}
Before the new one: say the Tamil for a promotion, then the Tamil for an ambition.
\end{quote}

\subsection*{You'll want to know: \ta{உழை}}
\textbf{\ta{உழை}} — \emph{uḻai} — \textquotedblleft{}to work hard; to toil\textquotedblright{}.

\begin{quote}
\ta{விவசாயிகள்} \ta{காலை} \ta{முதல்} \ta{மாலை} \ta{வரை} \ta{உழைக்கிறார்கள்}.
\end{quote}

\begin{quote}
\emph{vivacāyikaḷ kālai mutal mālai varai uḻaikkiṟārkaḷ.}
\end{quote}

\begin{quote}
Farmers toil from morning till evening.
\end{quote}

Past (I …): \ta{உழைத்தேன்}, \emph{uḻaittēṉ}.

A strong verb: \ta{உழைக்கிறேன்}, \ta{உழைத்தேன்}. The noun is \ta{உழைப்பு} (\emph{uḻaippu}), hard work.

\begin{cousinweb}
\textbf{\ta{உழை}} (\emph{uḻai}) is native Dravidian: DEDR 680 gives \ta{உழை} (\emph{uḻai}), \textquotedblleft{}to labour hard\textquotedblright{}, beside \ta{உழ} (\emph{uḻa}), \textquotedblleft{}to suffer, to toil\textquotedblright{}.
\end{cousinweb}

\begin{grammarlens}[title={What you've built}]
One more word for describing what you lived and felt.
\end{grammarlens}

\subsection*{Your turn}
\begin{itemize}
\raggedright
  \item \emph{Say it:} \emph{uḻai}
  \item \emph{Say it:} \emph{uḻai}, once more
  \item \emph{Say it:} \emph{vivacāyikaḷ kālai mutal mālai varai uḻaikkiṟārkaḷ.} — Farmers toil from morning till evening.
  \item \emph{Recall:} say the Tamil for a promotion, then the Tamil for an ambition, then say \emph{uḻai} again
\end{itemize}

\begin{tcolorbox}[breakable,colback=teal!4,colframe=teal!35!black,title={Before you move on}]
What does \ta{உழை} mean? (\textquotedblleft{}To work hard; to toil\textquotedblright{}.) Say it once more.
\end{tcolorbox}
\section[\texorpdfstring{\ta{ஊக்குவி} (ūkkuvi)}{ஊக்குவி (ūkkuvi)}]{\ta{ஊக்குவி} (ūkkuvi) — to encourage}

\label{lesson:TA-C509-ukkuvi}



\begin{quote}
Before the new one: say the Tamil for an ambition, then the Tamil for to work hard.
\end{quote}

\subsection*{You'll want to know: \ta{ஊக்குவி}}
\textbf{\ta{ஊக்குவி}} — \emph{ūkkuvi} — \textquotedblleft{}to encourage\textquotedblright{}.

\begin{quote}
\ta{ஆசிரியர்} \ta{எங்களைப்} \ta{பாடல்} \ta{போட்டியில்} \ta{பங்கேற்க} \ta{ஊக்குவித்தார்}.
\end{quote}

\begin{quote}
\emph{āciriyar eṅkaḷaip pāṭal pōṭṭiyil paṅkēṟka ūkkuvittār.}
\end{quote}

\begin{quote}
The teacher encouraged us to take part in the singing contest.
\end{quote}

Past (I …): \ta{ஊக்குவித்தேன்}, \emph{ūkkuvittēṉ}.

The person takes -\ta{ஐ}, the action is an infinitive: \ta{அவனை} \ta{வர} \ta{ஊக்குவி}, to encourage him to come.

\begin{cousinweb}
\textbf{\ta{ஊக்குவி}} (\emph{ūkkuvi}) is native Dravidian: DEDR 729 gives \ta{ஊக்கு} (\emph{ūkku}), \textquotedblleft{}to make an effort\textquotedblright{}, and \ta{ஊக்கம்} (\emph{ūkkam}), \textquotedblleft{}zeal, energy\textquotedblright{}. The ending -\ta{வி} makes a verb \textquotedblleft{}to give someone zeal\textquotedblright{}.
\end{cousinweb}

\begin{grammarlens}[title={What you've built}]
One more word for describing what you lived and felt.
\end{grammarlens}

\subsection*{Your turn}
\begin{itemize}
\raggedright
  \item \emph{Say it:} \emph{ūkkuvi}
  \item \emph{Say it:} \emph{ūkkuvi}, once more
  \item \emph{Say it:} \emph{āciriyar eṅkaḷaip pāṭal pōṭṭiyil paṅkēṟka ūkkuvittār.} — The teacher encouraged us to take part in the singing contest.
  \item \emph{Recall:} say the Tamil for an ambition, then the Tamil for to work hard, then say \emph{ūkkuvi} again
\end{itemize}

\begin{tcolorbox}[breakable,colback=teal!4,colframe=teal!35!black,title={Before you move on}]
What does \ta{ஊக்குவி} mean? (\textquotedblleft{}To encourage\textquotedblright{}.) Say it once more.
\end{tcolorbox}
\section[\texorpdfstring{\ta{நிபுணர்} (nipuṇar)}{நிபுணர் (nipuṇar)}]{\ta{நிபுணர்} (nipuṇar) — an expert}

\label{lesson:TA-C509-nipunar}



\begin{quote}
Before the new one: say the Tamil for to work hard, then the Tamil for to encourage.
\end{quote}

\subsection*{You'll want to know: \ta{நிபுணர்}}
\textbf{\ta{நிபுணர்}} — \emph{nipuṇar} — \textquotedblleft{}an expert\textquotedblright{}.

\begin{quote}
\ta{அவர்} \ta{கணினித்} \ta{துறையில்} \ta{ஒரு} \ta{நிபுணர்}.
\end{quote}

\begin{quote}
\emph{avar kaṇiṉit tuṟaiyil oru nipuṇar.}
\end{quote}

\begin{quote}
He is an expert in the computer field.
\end{quote}

Plural: \ta{நிபுணர்கள்}, \emph{nipuṇarkaḷ}.

A sentence of the shape \textquotedblleft{}X is a Y\textquotedblright{} needs no verb in Tamil: \ta{அவர்} \ta{ஒரு} \ta{நிபுணர்}.

\begin{cousinweb}
\textbf{\ta{நிபுணர்}} (\emph{nipuṇar}) is borrowed from Sanskrit \emph{nipuṇa}, \textquotedblleft{}clever, skilful\textquotedblright{} (Monier-Williams), with the Tamil personal ending -\ta{அர்}.
\end{cousinweb}

\begin{grammarlens}[title={What you've built}]
One more word for describing what you lived and felt.
\end{grammarlens}

\subsection*{Your turn}
\begin{itemize}
\raggedright
  \item \emph{Say it:} \emph{nipuṇar}
  \item \emph{Say it:} \emph{nipuṇar}, once more
  \item \emph{Say it:} \emph{avar kaṇiṉit tuṟaiyil oru nipuṇar.} — He is an expert in the computer field.
  \item \emph{Recall:} say the Tamil for to work hard, then the Tamil for to encourage, then say \emph{nipuṇar} again
\end{itemize}

\begin{tcolorbox}[breakable,colback=teal!4,colframe=teal!35!black,title={Before you move on}]
What does \ta{நிபுணர்} mean? (\textquotedblleft{}An expert\textquotedblright{}.) Say it once more.
\end{tcolorbox}
